\documentclass{article}
\usepackage{hyperref}
\usepackage[a4paper]{geometry}
\usepackage{fancyhdr}
\pagestyle{fancy}
\lhead{Intraindustrieller Handel}
\rhead{Februar 2026}
\begin{document}
\section{Intraindustrieller Handel}
Der \emph{intraindustrielle} Handel beschreibt Handel zwischen Nationen, welcher innerhalb (\emph{intra}) der gleichen Industrie stattfindet. Dieser ist das Gegenteil der interindustriellen Handels, bei welchem Nationen waren unterschiedlicher Arten handeln, den \hyperref[komparative Kostenvorteilen]{Komparative Kostenvorteile} nach welches am produktivsten ist.
 
\subsection{Erklärung} 
Wird sich nun der Handel von generelle Warengruppen zwischen zwei Industrienationen angeguckt, so fällt auf dass dies nicht mehr der Fall ist; Warengrupppen werden sowohl importiert als auch exportiert. Dies ist darin zu begründen, dass zwei Industrienationen in der Regel ähnlich produktiv sind; die komparativen Kostenvorteile verschwinden. Zudem können sie komplexere Produkte produzieren, welche für mehr relevante Varietät erlauben.
 
Es muss aber angemerkt werden, dass wenn sich der Handel genauer angeguckt wird, die Warengruppen spezifiziert werden, der intraindustrielle Handel verschwindet. 
 
\subsection{Grubel-Lloyd} 
Der intraindustrielle Handel kann mit dem \emph{Grubel-Lloyd-Index} gemessen werden. Wird eine Warengruppe importiert aber keineswegs exportiert oder exportiert und keineswegs importiert, so liegt dieser bei 0. Wird eine Warengruppe so viel importiert wie sie exportiert wird, liegt der Index bei 1. 
\end{document}